\documentclass[10pt,a4paper]{amsart}
\usepackage[margin=19mm]{geometry}
\usepackage{amsmath,amssymb,mathtools}
\usepackage[scr=rsfs,cal=euler]{mathalfa}
\usepackage{xfrac}
\usepackage[unicode,hidelinks,bookmarks=false]{hyperref}
\newcommand{\ie}{i.e.\ }
\newcommand{\cf}{cf.\ }
\newcommand{\resp}{respectively\ }
\newcommand{\Subsection}{\subsection}
\providecommand{\nobreakdash}{\nobreak\hbox{-}}
\setlength{\emergencystretch}{2em}
\multlinegap0pt
\usepackage{amssymb}
\usepackage{mathtools}
\newtheorem{lemma}{Lemma}
\newtheorem{proposition}{Proposition}
\newcommand{\R}{\mathbb R}
\newcommand{\TV}{\mathrm{TV}}
\newcommand{\Var}{\operatorname{Var}}
\newcommand{\be}{\beta}
\newcommand{\dd}{\,\mathrm d}
\newcommand{\ee}{\mathrm e}
\title{Spectral gaps for mean-field Coulomb gases}
\author{Simon Becker, Angeliki Menegaki}
\date{Source: arXiv:2608.21215v1 (CC BY 4.0)}
\begin{document}
\maketitle

\noindent\emph{Source and licence.} Spectral gaps for mean-field Coulomb gases. Version arXiv:2608.21215v1 (CC BY 4.0), distributed by arXiv under the Creative Commons Attribution 4.0 International licence, http://creativecommons.org/licenses/by/4.0/, as recorded in the arXiv licence metadata for this exact version and shown on the abstract page https://arxiv.org/abs/2608.21215v1. Full licence notice: \texttt{licenses/arxiv-2608.21215-CC-BY-4.0.txt}.\par\smallskip

\noindent\textit{Editorial scope.} Counted unit: the article's proposition prop:moving-center (source0.tex lines 1410--1436). The article's complete original proof of it is delivered with it (source0.tex lines 1438--1696, one \texttt{proof} environment of 259 lines). No mathematics is written, repaired or re-derived: the statement and the proof are the article's own lines, in the article's own notation, and no retained line is substituted or re-flowed.  Retained uncounted as the source-local context the proof needs: lines 478--483; lines 499--511.  Nothing was omitted from any retained span, so the package carries no omission record.  No retained line contains a citation, so the wrapper carries no bibliography and no bibliography entry.  No retained line contains a figure environment, an image-inclusion command or drawing code, so no image is delivered and the package declares no asset dependency.  Editorial formatting only, in addition: an AMS \texttt{amsart} preamble that restates the article's own declarations and macros used by the retained lines, namely the environments \texttt{lemma}, \texttt{proposition} on separate counters, together with the standard packages those lines need.  The retained environments print in this document as Lemma 1, Proposition 1 in order of appearance, whereas the article prints the same items with the numbering of its own full document.  The complete native source of the article as distributed in the arXiv e-print for this exact version is bundled unchanged as \texttt{sources/208233/source0.tex}.
\begin{lemma}\label{lem:gaussian-moment}
It holds that \begin{equation*} 
 M_{d,\be}:=
 \sup_{y\in\R^d}\int_{\R^d}|x-y|^{2-d}\dd\gamma_\be(x)<\infty.
\end{equation*}
\end{lemma}

\begin{proposition}
\label{prop:basic-onesite}
If \(\Gamma(\rho)\le(2M_{d,\be})^{-1}\), then
\[
 \begin{gathered}
 Z_{\rho}\ge\tfrac12,\quad
 \nu_{\rho}\le2\gamma_\be,\quad
 \sup_{y\in\R^d}\int_{\R^d}|x-y|^{2-d}\dd\nu_{\rho}(x)
 \le2M_{d,\be}, \ \|\nu_{\rho}-\gamma_\be\|_{\TV}
 \le4M_{d,\be}\Gamma(\rho).
 \end{gathered}
\]
\end{proposition}

\begin{proposition}
\label{prop:moving-center}
Let \(\rho\) be a finite positive point measure, and let \(a\ge0\) be the
charge of an additional Coulomb pole.  Assume that for $M_{d,\beta}$ the constant in Lemma \ref{lem:gaussian-moment}
\[
 \Gamma(\rho)+a\le\frac{1}{2M_{d,\be}}.
\]
For \(y\in\R^d\), define
\[
 \nu^y:=\nu_{\rho+a\delta_y}, \qquad\text{i.e. }
 \nu^y(dx)
 =
 \frac{\ee^{-a|x-y|^{2-d}}}{
   \int_{\R^d}\ee^{-a|u-y|^{2-d}}\dd\nu_\rho(u)}
 \nu_\rho(dx).
\]
Then, for every \(y\in\R^d\),
\begin{equation}\label{eq:add-remove-TV}
 \|\nu^y-\nu_\rho\|_{\TV}
 \le4M_{d,\be}a.
\end{equation}
Moreover, for every \(y,y'\in\R^d\),
\begin{equation}\label{eq:move-TV}
 \|\nu^y-\nu^{y'}\|_{\TV}
 \le8M_{d,\be}a.
\end{equation}
\end{proposition}

\begin{proof}
We first prove the add-remove estimate \eqref{eq:add-remove-TV} for a fixed
pole location \(y\).  The proof uses a bounded truncation of the singular
kernel, differentiates along the interpolation that turns on the pole, and
then removes the truncation by dominated convergence.

Fix \(y\in\R^d\).  For \(R>0\), define the bounded truncation
\[
 \phi_R(x):=\min\{|x-y|^{2-d},R\},
 \qquad x\in\R^d,
\]
where we set \(\phi_R(y)=R\).  For \(0\le t\le1\), set
\[
 Z_{t,R}
 :=
 \int_{\R^d}\ee^{-ta\phi_R(u)}\dd\nu_\rho(u) \text{ and }
 \nu_{t,R}(d x)
 :=
 Z_{t,R}^{-1}\ee^{-ta\phi_R(x)}\nu_\rho(d x).
\]
Since
\[
 \ee^{-taR}\le \ee^{-ta\phi_R(x)}\le1,
\]
one has \(0<Z_{t,R}\le1\), so \(\nu_{t,R}\) is a well-defined
probability measure.  Notice also that
\[
 \nu_{0,R}=\nu_\rho.
\]

We first obtain a bound on the \(\phi_R\)-moment of \(\nu_{t,R}\).
Because \(0\le\phi_R\le R\), differentiation under the integral is
justified by dominated convergence.  Thus
\begin{align*}
 \frac{\dd}{\dd t}Z_{t,R}
 &=-a\int_{\R^d}\phi_R(u)\ee^{-ta\phi_R(u)}
   \dd\nu_\rho(u)=-aZ_{t,R}\int_{\R^d}\phi_R(u)\dd\nu_{t,R}(u)
  =-aZ_{t,R}\nu_{t,R}(\phi_R).
\end{align*}

Define
\[
 N_{t,R}
 :=
 \int_{\R^d}\phi_R(u)\ee^{-ta\phi_R(u)}
 \dd\nu_\rho(u).
\]
Then $\nu_{t,R}(\phi_R)=\frac{N_{t,R}}{Z_{t,R}}.$ Again using the boundedness of \(\phi_R\), we may differentiate \(N_{t,R}\):
\begin{align*}
 \frac{\dd}{\dd t}N_{t,R}
 &=
 -a\int_{\R^d}\phi_R(u)^2\ee^{-ta\phi_R(u)}
 \dd\nu_\rho(u)=
 -aZ_{t,R}\nu_{t,R}(\phi_R^2).
\end{align*}
The quotient rule therefore gives
\begin{align*}
 \frac{\dd}{\dd t}\nu_{t,R}(\phi_R)
 &=
 \frac{N_{t,R}'Z_{t,R}-N_{t,R}Z_{t,R}'}
      {Z_{t,R}^2}=
 -a\nu_{t,R}(\phi_R^2)
 +a\nu_{t,R}(\phi_R)^2=
 -a\Var_{\nu_{t,R}}(\phi_R)
 \le0.
\end{align*}
Hence the map \(t\mapsto\nu_{t,R}(\phi_R)\) is nonincreasing.  Since
\(\nu_{0,R}=\nu_\rho\), it follows that
\[
 \nu_{t,R}(\phi_R)
 \le\nu_{0,R}(\phi_R)
 =\nu_\rho(\phi_R).
\]
Moreover,
\[
 0\le\phi_R(x)\le |x-y|^{2-d}
 \qquad\text{for }x\ne y.
\]
Since \(\nu_\rho\) is absolutely continuous with respect to Lebesgue
measure, the point \(y\) has zero \(\nu_\rho\)-measure.  Proposition
\ref{prop:basic-onesite} therefore gives
\begin{align*}
 \nu_{t,R}(\phi_R)
 &\le
 \int_{\R^d}\phi_R(x)\dd\nu_\rho(x)\le
 \int_{\R^d}|x-y|^{2-d}\dd\nu_\rho(x)\le 2M_{d,\be}.
\end{align*}
This estimate is uniform in \(t\in[0,1]\), \(R>0\), and \(y\in\R^d\). To estimate the change along the interpolation
\(t\mapsto\nu_{t,R}\), let \(h:\R^d\to\R\) be bounded and measurable.
Write
\[
 H_{t,R}
 :=
 \int_{\R^d}h(u)\ee^{-ta\phi_R(u)}\dd\nu_\rho(u),
\]
so that $\nu_{t,R}(h)=\frac{H_{t,R}}{Z_{t,R}}.$
Since \(h\) and \(\phi_R\) are bounded,
\begin{align*}
 \frac{\dd}{\dd t}H_{t,R}
 &=
 -a\int_{\R^d}h(u)\phi_R(u)
 \ee^{-ta\phi_R(u)}\dd\nu_\rho(u)=
 -aZ_{t,R}\nu_{t,R}(h\phi_R).
\end{align*}
Using the quotient rule and the preceding formula for \(Z_{t,R}'\), we
obtain
\begin{align*}
 \frac{\dd}{\dd t}\nu_{t,R}(h)
 &=
 \frac{H_{t,R}'Z_{t,R}-H_{t,R}Z_{t,R}'}
      {Z_{t,R}^2}=
 -a\nu_{t,R}(h\phi_R)
 +a\nu_{t,R}(h)\nu_{t,R}(\phi_R)=
 -a\operatorname{Cov}_{\nu_{t,R}}(h,\phi_R).
\end{align*}

Assume now that \(\|h\|_\infty\le1\).  Since \(\nu_{t,R}\) is a
probability measure $|\nu_{t,R}(h)|\le1$
and consequently
\[
 |h(x)-\nu_{t,R}(h)|\le2
 \qquad\text{for every }x\in\R^d.
\]
Because \(\phi_R\ge0\),
\begin{align*}
 \left|
 \operatorname{Cov}_{\nu_{t,R}}(h,\phi_R)
 \right|
 &=
 \left|
 \int_{\R^d}
 \bigl(h(x)-\nu_{t,R}(h)\bigr)\phi_R(x)
 \dd\nu_{t,R}(x)
 \right|\le
 \int_{\R^d}
 |h(x)-\nu_{t,R}(h)|\phi_R(x)
 \dd\nu_{t,R}(x)\\
 &\le
 2\int_{\R^d}\phi_R(x)\dd\nu_{t,R}(x)=
 2\nu_{t,R}(\phi_R)\le4M_{d,\be}.
\end{align*}
It follows that
\[
 \left|
 \frac{\dd}{\dd t}\nu_{t,R}(h)
 \right|
 \le4M_{d,\be}a.
\]
Integrating this estimate over \(t\in[0,1]\), and using
\(\nu_{0,R}=\nu_\rho\), gives
\begin{align*}
 |\nu_{1,R}(h)-\nu_\rho(h)|
 &=
 \left|
 \int_0^1
 \frac{\dd}{\dd t}\nu_{t,R}(h)\dd t
 \right|\le
 \int_0^1
 \left|
 \frac{\dd}{\dd t}\nu_{t,R}(h)
 \right|\dd t \le4M_{d,\be}a.
\end{align*}
Taking the supremum over all measurable \(h\) with
\(\|h\|_\infty\le1\) yields $\|\nu_{1,R}-\nu_\rho\|_{\TV}
 \le4M_{d,\be}a.$

It remains to remove the truncation.  Set $\phi(x):=|x-y|^{2-d}$ for $x\ne y.$
For every \(x\ne y\),
\[
 \phi_R(x)\uparrow\phi(x)
 \qquad\text{and hence}\qquad
 \ee^{-a\phi_R(x)}\longrightarrow\ee^{-a\phi(x)}
 \quad\text{as }R\to\infty.
\]
Moreover, $0\le\ee^{-a\phi_R(x)}\le1.$
Since \(\nu_\rho(\{y\})=0\), dominated convergence gives
\[
 Z_{1,R}
 =
 \int_{\R^d}\ee^{-a\phi_R(x)}\dd\nu_\rho(x)
 \longrightarrow
 Z_1
 :=
 \int_{\R^d}\ee^{-a|x-y|^{2-d}}\dd\nu_\rho(x).
\]
The limiting normalizing constant is strictly positive because the
integrand is strictly positive on \(\R^d\setminus\{y\}\), a set of full
\(\nu_\rho\)-measure.  Thus \(Z_1>0\).

The densities of \(\nu_{1,R}\) and \(\nu^y\) with respect to \(\nu_\rho\)
are
\[
 p_R(x):=Z_{1,R}^{-1}\ee^{-a\phi_R(x)},
 \qquad
 p(x):=Z_1^{-1}\ee^{-a|x-y|^{2-d}},
\]
respectively.  We have
\begin{align*}
 \int_{\R^d}|p_R(x)-p(x)|\dd\nu_\rho(x)
 &\le
 \left|Z_{1,R}^{-1}-Z_1^{-1}\right|
 \int_{\R^d}\ee^{-a\phi_R(x)}\dd\nu_\rho(x)+
 Z_1^{-1}
 \int_{\R^d}
 \left|
 \ee^{-a\phi_R(x)}
 -\ee^{-a|x-y|^{2-d}}
 \right|
 \dd\nu_\rho(x)\\
 &=
 Z_{1,R}\left|Z_{1,R}^{-1}-Z_1^{-1}\right|+Z_1^{-1}
 \int_{\R^d}
 \left|
 \ee^{-a\phi_R(x)}
 -\ee^{-a|x-y|^{2-d}}
 \right|
 \dd\nu_\rho(x).
\end{align*}
The first term converges to zero because \(Z_{1,R}\to Z_1>0\), and the
second converges to zero by dominated convergence.  Hence
\[
 \int_{\R^d}|p_R(x)-p(x)|\dd\nu_\rho(x)\longrightarrow0.
\]
Since both measures are absolutely continuous with respect to \(\nu_\rho\),
with densities \(p_R\) and \(p\), respectively, the definition of the full
total-variation norm gives
\begin{align*}
 \|\nu_{1,R}-\nu^y\|_{\TV}
 &=
 \sup_{\|h\|_\infty\le1}
 \left|
 \int_{\R^d}h(x)\bigl(p_R(x)-p(x)\bigr)\dd\nu_\rho(x)
 \right|\le
 \int_{\R^d}|p_R(x)-p(x)|\dd\nu_\rho(x).
\end{align*}
The \(L^1(\nu_\rho)\)-convergence of the densities established above now
implies
\[
 \|\nu_{1,R}-\nu^y\|_{\TV}\longrightarrow0.
\]
Therefore,  $\|\nu^y-\nu_\rho\|_{\TV}
 \le
 \|\nu^y-\nu_{1,R}\|_{\TV}
 +\|\nu_{1,R}-\nu_\rho\|_{\TV}.$
Letting \(R\to\infty\) gives $\|\nu^y-\nu_\rho\|_{\TV}
 \le4M_{d,\be}a,$
which proves \eqref{eq:add-remove-TV}.

Finally, applying \eqref{eq:add-remove-TV} separately to the poles at
\(y\) and \(y'\), and using the triangle inequality, gives
\begin{align*}
 \|\nu^y-\nu^{y'}\|_{\TV}
 &\le
 \|\nu^y-\nu_\rho\|_{\TV}
 +\|\nu_\rho-\nu^{y'}\|_{\TV}\le
 4M_{d,\be}a+4M_{d,\be}a=8M_{d,\be}a.
\end{align*}
This proves \eqref{eq:move-TV}.
\end{proof}

\end{document}
